\documentclass[a4paper,11pt]{article}
\usepackage[T1]{fontenc}
\usepackage[utf8]{inputenc}
\usepackage{lmodern}
\usepackage[top=2cm,bottom=2cm,left=2cm,right=2cm]{geometry}
\usepackage{booktabs}
\usepackage{parskip}
\usepackage{microtype}
\usepackage{hyperref}
\hypersetup{colorlinks=true,linkcolor=black,urlcolor=blue,citecolor=black}

\title{\textbf{KKO 2025/2026 -- Grayscale Image Compressor}\\[0.3em]
       \large LZSS-Based Dictionary Compression}
\author{Martin Pribylina (xpriby19)}
\date{May 2026}

\begin{document}
\maketitle
\thispagestyle{empty}

% -----------------------------------------------------------------------
\section{Problem Overview}
% -----------------------------------------------------------------------

The task is to write a console application in C++ that compresses 8-bit
grayscale images stored in the RAW format using a lossless
\textbf{data compression algorithm} based on the sliding-window principle.
The program \texttt{lz\_codec} has two operating modes: compression
(\texttt{-c}) and decompression (\texttt{-d}).  During compression the user
can optionally enable a pixel-difference pre-processing step (\texttt{-m})
and an adaptive per-block scanning mode (\texttt{-a}).  The
decompressor needs to reconstruct the original image must be embedded in the
compressed file so that no extra parameters are required when decompressing.

% -----------------------------------------------------------------------
\section{Chosen Solution}
% -----------------------------------------------------------------------

\subsection*{Algorithm: LZSS}

LZSS (Lempel--Ziv--Storer--Szymanski, 1982) is a sliding-window dictionary
method.  For each input position it searches a look-back buffer for the
longest repetition.  If a match of length $\geq$ \texttt{MIN\_MATCH} is
found, a \emph{match token} \texttt{(offset, length)} is emitted; otherwise
a \emph{literal byte} is emitted.  A single \textbf{flag bit} distinguishes
the two cases, so short literals are not penalised with unused fields (unlike
plain LZ77).

Parameters used in this implementation:

\begin{center}
\begin{tabular}{lll}
\toprule
Parameter & Value & Reason \\
\midrule
Search window & 4\,096\,B & 12-bit offset, covers one full image row and beyond \\
Lookahead & 32\,B & 5-bit length field, maximum match 34 bytes \\
Min match & 3\,B & A match token (18 bits) beats three literals (27 bits) \\
Flag packing & 8 flags/byte & Keeps the bit-stream byte-aligned \\
\bottomrule
\end{tabular}
\end{center}

The encoder uses \textbf{hash chains} (2-byte hash, chain depth 128) to
locate matches quickly.  The decoder only reads and copies data, with no
search at all.

\subsection*{Scanning and Serialization}

Because images are two-dimensional, the pixels must be arranged into a
one-dimensional byte sequence before compression.  The order of this
sequence has a direct effect on how many repeating patterns LZSS can find.

\textbf{Static mode} (default): pixels are scanned row by row (horizontal).
The entire image is treated as one contiguous byte stream.

\textbf{Adaptive mode} (\texttt{-a}): the image is divided into
$16\!\times\!16$ pixel blocks.  For each block the \emph{sum of absolute
differences} (SAD) is computed along rows and along columns; the direction
with the lower SAD is chosen because smoother transitions produce more
repetition and longer matches.  If LZSS cannot compress a block (encoded
size $\geq$ raw size), the block is stored as-is.  Two metadata bits per
block (scan direction and compression flag) are packed into a compact bit
array written before the block data.

\subsection*{Pixel-Difference Model (\texttt{-m})}

Before LZSS encoding, each byte is replaced by its difference from the
previous byte:
\[
  \delta_i = x_i - x_{i-1} \pmod{256}, \quad \delta_0 = x_0.
\]
For smooth images, differences cluster near zero, creating more repeating
byte sequences and therefore longer LZSS matches.  The inverse transform
($x_i = x_{i-1} + \delta_i$) is applied after decoding.

\subsection*{File Format}

The compressed file header (9~bytes, little-endian):

\begin{center}
\begin{tabular}{cll}
\toprule
Offset & Size & Field \\
\midrule
0 & 2\,B & \texttt{uint16\_t} image width \\
2 & 2\,B & \texttt{uint16\_t} image height \\
4 & 1\,B & flags: bit\,0 = model used, bit\,1 = adaptive, bit\,2 = raw fallback \\
5 & 4\,B & \texttt{uint32\_t} original pixel count \\
\bottomrule
\end{tabular}
\end{center}

Bit\,2 of flags (\emph{raw fallback}) is set when compression would produce
output larger than the input.  In that case the payload is the original raw
pixels and the output size equals the input size plus the 9-byte header.

% -----------------------------------------------------------------------
\section{Key Data Structures}
% -----------------------------------------------------------------------

\textbf{BitWriter / BitReader} (\texttt{bitio.hpp}): accumulate bits into
bytes using a 1-byte buffer.  \texttt{flush()} zero-pads the final byte.
Reading and writing arbitrary bit widths allows compact LZSS tokens without
byte-alignment waste.

\textbf{Window} (circular buffer, 4096 bytes): maintained identically by
both encoder and decoder so the dictionary does not need to be transmitted.

\textbf{HashChains}: two arrays (\texttt{head[65536]}, \texttt{prev[4096]})
index the window by the 2-byte hash of each position.  Walking the chain up
to depth 128 finds long matches without scanning the entire window.

\textbf{FlagGroup}: accumulates up to 8 tokens (literals or match
references), then writes one flag byte followed by the payloads in order.
This keeps the bit-stream layout identical on both sides.

\textbf{Block / BlockMeta}: holds the $16\!\times\!16$ pixel sub-image, its
chosen scan direction (0 = horizontal, 1 = vertical), and its compression
flag.  Block metadata is bit-packed (2 bits per block) into a compact array
before the block payloads.

% -----------------------------------------------------------------------
\section{Results}
% -----------------------------------------------------------------------

All measurements are on 512$\times$512 (262\,144\,B) 8-bit grayscale images
compiled with \texttt{g++ -O2 -std=c++17} on a Linux x86-64 host.
Entropy $H$ is the first-order (symbol-independent) Shannon entropy of each
image.

\begin{center}
\small
\setlength{\tabcolsep}{4pt}
\begin{tabular}{lrrrrrrrrrr}
\toprule
 & \multicolumn{2}{c}{Static} & \multicolumn{2}{c}{Static+M} &
   \multicolumn{2}{c}{Adaptive} & \multicolumn{2}{c}{Adaptive+M} & \\
\cmidrule(lr){2-3}\cmidrule(lr){4-5}\cmidrule(lr){6-7}\cmidrule(lr){8-9}
Image & bpp & ms & bpp & ms & bpp & ms & bpp & ms & $H$ \\
\midrule
cb      & 0.56 &  47 & 0.57 &  44 & 0.76 & 203 & 0.77 & 197 & 1.00 \\
cb2     & 0.68 &  49 & 0.70 &  51 & 2.29 & 200 & 0.91 & 195 & 6.91 \\
df1h    & 0.57 &  44 & 0.56 &  43 & 1.82 & 199 & 0.91 & 201 & 8.00 \\
df1hvx  & 0.61 &  57 & 0.67 &  61 & 2.06 & 198 & 1.65 & 208 & 4.51 \\
df1v    & 0.58 &  50 & 0.56 &  42 & 1.82 & 209 & 0.91 & 188 & 8.00 \\
nk01    & 8.00 & 180 & 8.00 & 177 & 8.00 & 151 & 8.00 & 152 & 6.47 \\
shp     & 0.63 &  48 & 0.62 &  46 & 2.10 & 212 & 2.20 & 215 & 0.87 \\
shp1    & 1.25 &  66 & 1.29 &  67 & 2.02 & 210 & 2.10 & 211 & 1.90 \\
shp2    & 1.46 &  79 & 1.54 &  83 & 1.92 & 205 & 2.01 & 208 & 1.87 \\
\midrule
\textbf{avg} & \textbf{1.59} & \textbf{69} & \textbf{1.61} & \textbf{68} & \textbf{2.53} & \textbf{198} & \textbf{2.16} & \textbf{197} & \textbf{4.39} \\
\bottomrule
\end{tabular}
\end{center}

\smallskip
\textbf{Observations:}

\begin{itemize}\setlength\itemsep{2pt}
  \item \textbf{Static mode is the best overall (1.59 bpp average)}, which
        is perhaps surprising.  Adaptive mode breaks the image into small
        independent blocks and each block starts with an empty LZSS window,
        losing all the match history from previous blocks.  Static mode
        compresses the whole image as one stream, so the sliding window
        accumulates repetitions across entire rows and even across rows,
        giving it a significant advantage on most images.

  \item \textbf{The pixel-difference model helps only for correlated images.}
        For gradients and smooth images (df1h, df1v) the model drops the
        static result from 0.57 to 0.56 bpp and the adaptive result from
        1.82 to 0.91 bpp -- a major win.  For images with sharp edges and
        little spatial correlation (shp1, shp2) the model slightly hurts,
        because pixel differences there are just as random as the original
        values.  This was expected from the theory.

  \item \textbf{Adaptive~+~Model recovers most of the adaptive penalty} for
        correlated images (cb2: 2.29~$\to$~0.91, df1h: 1.82~$\to$~0.91),
        since the delta transform effectively re-creates long repeating runs
        even in short 256-byte blocks.

  \item \textbf{nk01 cannot be compressed.}  Its entropy is 6.47 bpp and
        it contains no repeating structure.  All four modes produce output
        equal to the input (8.00 bpp includes the 9-byte header overhead)
        because the raw fallback is triggered.  This was expected.

  \item \textbf{Adaptive mode is about four times slower} than static mode.
        This is expected because it runs two passes per block: first the SAD
        heuristic to decide the scan direction, then the actual LZSS
        encoding.  Static mode is one single pass over the whole image.
\end{itemize}

% -----------------------------------------------------------------------
\section*{References}
% -----------------------------------------------------------------------

\begin{enumerate}\setlength\itemsep{1pt}
  \item Přednášky k předmětu Kódování a komprese dat, FIT VUT (2025/2026).
  \item Storer, J.A., Szymanski, T.G.: \emph{Data compression via textual substitution},
        JACM, 1982.
  \item Salomon, D.: \emph{Data Compression -- The Complete Reference},
        Springer, 2000.
\end{enumerate}

\end{document}
